% ═══════════════════════════════════════════════════════════════
% LEHRERHINWEISE - Physik Klasse 6: Bewegungsarten erkennen
% Stunde 1 von 4 | Reihe: Bewegungen
% ═══════════════════════════════════════════════════════════════

\documentclass[11pt, a4paper]{article}

\newif\ifswmodus
\swmodusfalse

\usepackage[utf8]{inputenc}
\usepackage[T1]{fontenc}
\usepackage[ngerman]{babel}

\usepackage{helvet}
\renewcommand{\familydefault}{\sfdefault}

\usepackage[margin=2cm, top=2cm, bottom=2cm]{geometry}
\usepackage{enumitem}
\usepackage{tabularx}
\usepackage{booktabs}
\usepackage{array}
\usepackage{longtable}
\usepackage{amssymb}

\usepackage{xcolor}
\usepackage{tcolorbox}
\tcbuselibrary{skins, breakable}

\usepackage{fancyhdr}
\usepackage{lastpage}

\setlist{nosep, leftmargin=1.5em}

% ═══════════════════════════════════════════════════════════════
% FARBEN
% ═══════════════════════════════════════════════════════════════

\definecolor{physik-color}{HTML}{2c5aa0}
\definecolor{hellgrau-color}{HTML}{F0F0F0}
\definecolor{grau-color}{HTML}{666666}
\definecolor{success-color}{HTML}{2a7a4b}
\definecolor{warning-color}{HTML}{b35c00}
\definecolor{error-color}{HTML}{c62828}
\definecolor{basis-color}{HTML}{2a7a4b}
\definecolor{standard-color}{HTML}{2c5aa0}
\definecolor{erweiterung-color}{HTML}{7b1fa2}

\colorlet{fachfarbe}{physik-color}
\colorlet{hellgrau}{hellgrau-color}
\colorlet{grau}{grau-color}
\colorlet{success}{success-color}
\colorlet{warning}{warning-color}
\colorlet{error}{error-color}
\colorlet{basis}{basis-color}
\colorlet{standard}{standard-color}
\colorlet{erweiterung}{erweiterung-color}

% ═══════════════════════════════════════════════════════════════
% VARIABLEN
% ═══════════════════════════════════════════════════════════════

\newcommand{\fach}{Physik}
\newcommand{\thema}{Bewegungsarten erkennen}
\newcommand{\klasse}{6}
\newcommand{\stunde}{01}
\newcommand{\stundenanzahl}{4}
\newcommand{\reihe}{Bewegungen}
\newcommand{\dauer}{45 Min.}
\newcommand{\lerngruppengroesse}{24}

% ═══════════════════════════════════════════════════════════════
% KOPF-/FUSSZEILE
% ═══════════════════════════════════════════════════════════════

\pagestyle{fancy}
\fancyhf{}
\renewcommand{\headrulewidth}{0.4pt}
\lhead{\small\textcolor{fachfarbe}{\textbf{\fach{} Kl. \klasse}} -- \thema{} -- \textbf{Lehrerhinweise}}
\rhead{\small LH-\stunde}
\cfoot{\small\thepage/\pageref{LastPage}}

% ═══════════════════════════════════════════════════════════════
% BEFEHLE
% ═══════════════════════════════════════════════════════════════

\newcommand{\sternEins}{\textcolor{basis}{\textbf{(*)}}}
\newcommand{\sternZwei}{\textcolor{standard}{\textbf{(**)}}}
\newcommand{\sternDrei}{\textcolor{erweiterung}{\textbf{(***)}}}
\newcommand{\lsg}[1]{\textcolor{error}{\textbf{#1}}}

% ═══════════════════════════════════════════════════════════════
% BOXEN
% ═══════════════════════════════════════════════════════════════

\newtcolorbox{infobox}[1][]{
    colback=hellgrau,
    colframe=fachfarbe,
    boxrule=1pt,
    arc=0pt,
    left=8pt, right=8pt, top=6pt, bottom=6pt,
    fonttitle=\bfseries,
    title=#1,
    breakable
}

\newtcolorbox{warnbox}[1][]{
    colback=white,
    colframe=warning,
    boxrule=1pt,
    arc=0pt,
    left=8pt, right=8pt, top=6pt, bottom=6pt,
    fonttitle=\bfseries,
    title=#1,
    breakable
}

\newtcolorbox{tippbox}{
    colback=white,
    colframe=success,
    boxrule=1pt,
    arc=0pt,
    left=8pt, right=8pt, top=6pt, bottom=6pt,
    fonttitle=\bfseries,
    title=Tipp,
    breakable
}

% ═══════════════════════════════════════════════════════════════
% DOKUMENT
% ═══════════════════════════════════════════════════════════════

\begin{document}

\begin{center}
{\Large\textbf{\textcolor{fachfarbe}{Lehrerhinweise: \thema}}}\\[0.3em]
{\normalsize \fach{} Klasse \klasse{} -- Stunde \stunde{} von \stundenanzahl{} -- Reihe ``\reihe''}
\end{center}

\vspace{1em}

% ═══════════════════════════════════════════════════════════════
% 1. KURZÜBERSICHT
% ═══════════════════════════════════════════════════════════════

\section*{1. Kurzübersicht}

\begin{infobox}
\begin{tabular}{ll}
\textbf{Thema:} & Bewegungsarten erkennen \\
\textbf{Klasse:} & 6 (Orientierungsstufe, heterogen) \\
\textbf{Zeit:} & 45 Min. \\
\textbf{Lernziele:} & FZ1: Definition von Bewegung nennen (AFB I) \\
                    & FZ2: Zwei Bahnformen nennen (AFB I) \\
                    & FZ3: Zwei Bewegungsarten beschreiben (AFB I) \\
                    & FZ4: Beispiele zuordnen (AFB II) \\
\textbf{Methoden:} & Stillarbeit (LB), digitaler Test (MT) \\
\textbf{Material:} & LB-01, MT-01 (Tablets), Bild \\
\end{tabular}
\end{infobox}

% ═══════════════════════════════════════════════════════════════
% 2. STUNDENVERLAUF
% ═══════════════════════════════════════════════════════════════

\section*{2. Stundenverlauf}

\begin{longtable}{|c|p{2cm}|p{6.5cm}|c|p{2cm}|}
\hline
\textbf{Min} & \textbf{Phase} & \textbf{Inhalt} & \textbf{SF} & \textbf{Material} \\
\hline
\endhead

0--3 & Einstieg & Impuls: ``Welche Bewegungen seht ihr?'' Bild mit Auto, Karussell, Schlange zeigen. & PL & Bild \\
\hline

3--18 & Erarbeitung & SuS lesen LB selbstständig, füllen Kästen 1--3 aus. L geht herum, gibt Impulse. & EA & LB-01 \\
\hline

18--20 & Sicherung & Kurzes Abgleichen der Merksätze im PL. ``Wer hat geradlinig?'' & PL & --- \\
\hline

20--43 & Test & MT-01 digital bearbeiten. SuS wählen Niveau (*/\!**/\!***). & EA & MT-01, Tablets \\
\hline

43--45 & Abschluss & Tablets einsammeln, Ausblick: ``Nächste Stunde: Wie schnell?'' & PL & --- \\
\hline
\end{longtable}

% ═══════════════════════════════════════════════════════════════
% 3. DIFFERENZIERUNG
% ═══════════════════════════════════════════════════════════════

\section*{3. Differenzierung}

\begin{infobox}[Legende]
\sternEins{} Basis (BR) | \sternZwei{} Standard (MR) | \sternDrei{} Erweitert (GY)
\end{infobox}

\vspace{0.5em}

\begin{tabular}{|l|p{3.5cm}|p{3.5cm}|p{4cm}|}
\hline
& \sternEins{} \textbf{Basis} & \sternZwei{} \textbf{Standard} & \sternDrei{} \textbf{Erweitert} \\
\hline
\textbf{Ziel} & Begriffe zuordnen & + weitere Beispiele & + schwierigere Fälle \\
\hline
\textbf{MT-Fragen} & Q1--Q4 (25 BE) & Q1--Q5 (31 BE) & Q1--Q6 (35 BE) \\
\hline
\textbf{Hilfen} & Beispiele im LB & LB als Referenz & Nur Impulse \\
\hline
\end{tabular}

\vspace{0.5em}

\begin{tippbox}
Niveau im MT wird von SuS selbst gewählt. Bei Unsicherheit: ``Starte mit (*) und probiere (**) beim nächsten Mal.''
\end{tippbox}

% ═══════════════════════════════════════════════════════════════
% 4. HÄUFIGE SCHÜLERFEHLER
% ═══════════════════════════════════════════════════════════════

\section*{4. Häufige Schülerfehler}

\begin{warnbox}[Typische Fehler]
\begin{tabular}{|p{4cm}|p{8.5cm}|}
\hline
\textbf{Fehler} & \textbf{Reaktion / Hilfe} \\
\hline
\textbf{Kreisförmig $\neq$ krummlinig} & ``Ist ein Kreis eine gerade oder gekrümmte Linie?'' \\
\hline
\textbf{Gleichförmig = gleiche Form} & ``Worum geht es bei \textit{Bewegungs}art? Um Form oder Geschwindigkeit?'' \\
\hline
\textbf{Schaukel = kreisförmig} & ``Fährt die Schaukel einen ganzen Kreis?'' (Nein $\rightarrow$ krummlinig) \\
\hline
\textbf{Fußball = geradlinig} & ``Fliegt ein Fußball wirklich gerade durch die Luft?'' \\
\hline
\end{tabular}
\end{warnbox}

% ═══════════════════════════════════════════════════════════════
% 5. MATERIAL-CHECKLISTE
% ═══════════════════════════════════════════════════════════════

\section*{5. Material-Checkliste}

\begin{itemize}
    \item[$\Box$] LB-01.pdf kopiert (\lerngruppengroesse{} Exemplare)
    \item[$\Box$] MT-01.html auf Tablets getestet
    \item[$\Box$] Einstiegsbild (Auto, Karussell, Schlange)
    \item[$\Box$] WLAN funktioniert
    \item[$\Box$] Reserve-Tablets geladen
\end{itemize}

% ═══════════════════════════════════════════════════════════════
% 6. LÖSUNGEN LB-01
% ═══════════════════════════════════════════════════════════════

\section*{6. Lösungen LB-01}

\textbf{Kasten 1 -- Bewegung:}
\begin{itemize}
    \item Ein Körper ist in Bewegung, wenn er seinen \lsg{Ort} verändert.
    \item Die Bahnform beschreibt, wie der \lsg{Weg} aussieht.
    \item Die Bewegungsart beschreibt, ob sich die \lsg{Geschwindigkeit} ändert.
\end{itemize}

\textbf{Kasten 2 -- Bahnformen:}
\begin{itemize}
    \item Geradlinig: Der Körper bewegt sich auf einer \lsg{geraden Linie}.
    \item Krummlinig: Der Körper bewegt sich auf einer \lsg{gekrümmten} Bahn.
    \item Sonderfall kreisförmig: Der Körper bewegt sich auf einem \lsg{Kreis}.
\end{itemize}

\textbf{Kasten 3 -- Bewegungsarten:}
\begin{itemize}
    \item Gleichförmig: Die Geschwindigkeit bleibt \lsg{gleich / konstant}.
    \item Ungleichförmig: Die Geschwindigkeit \lsg{ändert} sich.
\end{itemize}

\textbf{Übung 5 -- Zuordnung:}
\begin{itemize}
    \item a) Rolltreppe: \lsg{geradlinig, gleichförmig}
    \item b) Schaukel: \lsg{krummlinig, ungleichförmig}
    \item c) Aufzug (volle Fahrt): \lsg{geradlinig, gleichförmig}
    \item d) Fußball beim Schuss: \lsg{krummlinig, ungleichförmig}
\end{itemize}

% ═══════════════════════════════════════════════════════════════
% 7. MT-LÖSUNGEN (KURZFORM)
% ═══════════════════════════════════════════════════════════════

\section*{7. MT-Lösungen (Kurzform)}

\begin{tabular}{|l|p{7cm}|c|}
\hline
\textbf{Nr} & \textbf{Lösung} & \textbf{BE} \\
\hline
Q1 & Was ist Bewegung? Ort, Bahnform, Weg, Bewegungsart, Geschwindigkeit & 5 \\
\hline
Q2 & Bahnformen: geradlinig, krummlinig & 2 \\
\hline
Q3 & Bewegungsarten: gleichförmig, ungleichförmig & 2 \\
\hline
Q4 & Zuordnung (8 Beispiele $\times$ 2 Kategorien) & 16 \\
\hline
\textbf{(*) Summe} & & \textbf{25 BE} \\
\hline
Q5 & Zuordnung (**): 3 weitere Beispiele $\times$ 2 & 6 \\
\hline
\textbf{(**) Summe} & & \textbf{31 BE} \\
\hline
Q6 & Zuordnung (***): 2 schwierige Fälle $\times$ 2 & 4 \\
\hline
\textbf{(***) Summe} & & \textbf{35 BE} \\
\hline
\end{tabular}

\vspace{0.5em}

\textbf{Q4 Lösungen (*)}:
\begin{itemize}[itemsep=1pt]
    \item Fallschirmspringer (konst. v): geradlinig, gleichförmig
    \item Ventilator-Blatt: krummlinig, gleichförmig
    \item Auto bremst: geradlinig, ungleichförmig
    \item Achterbahn im Looping: krummlinig, ungleichförmig
    \item Schnecke beim Kriechen: krummlinig, ungleichförmig
    \item Aufziehauto (wird langsamer): geradlinig, ungleichförmig
    \item Minutenzeiger einer Uhr: krummlinig, gleichförmig
    \item Stein fällt vom Baum: geradlinig, ungleichförmig
\end{itemize}

\textbf{Q5 Lösungen (**)}:
\begin{itemize}[itemsep=1pt]
    \item Mond um die Erde: krummlinig, gleichförmig
    \item Bobfahrer in der Kurve: krummlinig, ungleichförmig
    \item Skateboard bergab: geradlinig, ungleichförmig
\end{itemize}

\textbf{Q6 Lösungen (***)}:
\begin{itemize}[itemsep=1pt]
    \item Schmetterling fliegt: krummlinig, ungleichförmig
    \item Pendel einer Standuhr: krummlinig, ungleichförmig
\end{itemize}

% ═══════════════════════════════════════════════════════════════
% 8. AUSBLICK FOLGESTUNDEN
% ═══════════════════════════════════════════════════════════════

\section*{8. Ausblick: Folgestunden}

\begin{tabular}{|c|l|l|}
\hline
\textbf{Std.} & \textbf{Thema} & \textbf{Schwerpunkt} \\
\hline
1 & Bewegungsarten erkennen (diese Stunde) & Bahnform, Bewegungsart (qualitativ) \\
\hline
2 & Geschwindigkeit & v = s/t, km/h $\leftrightarrow$ m/s \\
\hline
3 & s(t)-Diagramme & Diagramme lesen und zeichnen \\
\hline
4 & Übung und Test & Wiederholung, Lernzielkontrolle \\
\hline
\end{tabular}

% ═══════════════════════════════════════════════════════════════
% 9. REFLEXIONSNOTIZEN
% ═══════════════════════════════════════════════════════════════

\section*{9. Reflexionsnotizen (nach Durchführung)}

\begin{tabular}{|p{7cm}|p{7cm}|}
\hline
\textbf{Was hat gut funktioniert?} & \textbf{Was muss angepasst werden?} \\
\hline
& \\[2.5cm]
\hline
\end{tabular}

\vspace{0.5em}

\textbf{Zeiteinhaltung:} $\Box$ okay \quad $\Box$ zu knapp \quad $\Box$ Zeit übrig

\textbf{Differenzierung:} $\Box$ okay \quad $\Box$ anpassen (wie: \rule{4cm}{0.4pt})

\textbf{MT:} $\Box$ funktioniert \quad $\Box$ technische Probleme

\end{document}
